% Appendix — extended inference benchmarks.
% Figures generated by benchmarks/plot_appendix.py from
% results/inference_sweep.csv and results/paradigm_bench_full.csv;
% PDFs in results/appendix_figs/. Needs: graphicx.
% All figures are single-column (\columnwidth) by construction, so this
% section is designed to fit within ~2 pages of a two-column layout.

\section{Extended Inference Benchmarks}
\label{sec:appendix_benchmarks}

We complement Section~\ref{sec:eval:inference} with finer-grained
sweeps: attention-variant characterization in isolation
(Figures~\ref{fig:appx_attention_type}--\ref{fig:appx_context}), and
paradigm scaling with generation length and batch size
(Figure~\ref{fig:appx_scaling}, shown for MHA; GQA and MLA are
qualitatively identical).

Figure~\ref{fig:appx_attention_type} isolates the attention variant on
an 8M backbone (B{=}1, fp32): GQA matches MHA at every compression
ratio, while MLA trades ${\approx}20\%$ decode throughput for its
compressed cache. Figure~\ref{fig:appx_batch} scales batch size: MHA
and GQA track ideal linear throughput scaling to B{=}128, MLA scales
linearly at a ${\approx}20\%$ lower slope, and prefill latency stays
batch-insensitive until B{=}128. Figure~\ref{fig:appx_context} scales
context length: MHA's KV cache grows ${\approx}4\times$ faster than
GQA-1/4's and ${\approx}6\times$ faster than MLA's --- the memory
saving that motivates both variants --- while decode throughput
remains context-insensitive in this range.

Figure~\ref{fig:appx_scaling} (left) scales generation length (medium
backbone, B{=}4, $S{=}32$ denoising steps): AR latency grows linearly
with generated length, while both diffusion paradigms finish in a
fixed $S$ passes, so their latency stays flat and throughput grows
linearly instead. The same figure (right) scales batch size on the
same paradigms (medium ${\approx}6.6$M backbone, G{=}128, fp32): at
this scale every decode step is dispatch- rather than compute-bound,
so diffusion's $S{=}32$ parallel passes beat AR's 128 sequential
steps at every batch size shown, and the AR crossover of
Figure~\ref{fig:inference_efficiency} sits beyond this range. MFU
confirms the regime: ${\approx}11\%$ for diffusion vs.\ ${<}1\%$ for
AR.

\begin{figure}[t]
\centering
\includegraphics[width=\columnwidth]{appx_attention_type.pdf}
\caption{\footnotesize Prefill latency and decode throughput per
  attention variant.}
\label{fig:appx_attention_type}
\end{figure}

\begin{figure}[t]
\centering
\includegraphics[width=\columnwidth]{appx_batch_size.pdf}
\caption{\footnotesize Decode throughput and prefill latency vs.\
  batch size, per attention variant.}
\label{fig:appx_batch}
\end{figure}

\begin{figure}[t]
\centering
\includegraphics[width=\columnwidth]{appx_context_len.pdf}
\caption{\footnotesize KV-cache memory and decode throughput vs.\
  context length, per attention variant.}
\label{fig:appx_context}
\end{figure}

\begin{figure}[t]
\centering
\includegraphics[width=\columnwidth]{appx_paradigm_scaling.pdf}
\caption{\footnotesize Latency, throughput, and memory/MFU vs.\
  generation length (left) and batch size (right), per paradigm.}
\label{fig:appx_scaling}
\end{figure}
